\begin{exercisechunk}
\item \textbf{Every string map has a next-token representation.}
\label[exercise]{ex:l16-autoregressive-representation}
\exerciseprofile{2}{1}{1}{\exproof\quad\exconstruction}{%
Separate the expressiveness of autoregressive output from the difficulty
of representing its next-token rule.}
Let $\Delta$ be a finite data alphabet and
$\Sigma=\Delta\cup\{\mathtt{SEP},\mathtt{EOS}\}$, where the two added
tokens are distinct and do not belong to $\Delta$.
Given $p:\Sigma^*\to\Sigma$, the runtime $G$ starts from
$h_0=x\mathtt{SEP}$, repeatedly appends $p(h)$ to the current prefix,
and stops when $p(h)=\mathtt{EOS}$, without appending the stop token.
It returns only the generated answer, excluding $x\mathtt{SEP}$.

For an arbitrary function $F:\Delta^*\to\Delta^*$, construct a total
next-token rule $p$ and prove that $G^p(x)=F(x)$ for every $x\in\Delta^*$.
Your construction must define $p$ on all strings in $\Sigma^*$ and cover
the case $F(x)$ is empty. Neither $F$ nor $p$ is assumed computable.
\end{exercisechunk}
